\chapter{Einleitung} \label{ch:Einleitung}

\section{Aufgabe und Motivation}
\label{sec:AuM}

Ziel dieses Versuchs ist die Bestimmung der Halbwertszeiten radioaktiver Isotope durch Aktivierung stabiler Kernspektren mittels thermischer Neutronen. Insbesondere sollen die Zerfallskonstanten und damit verbunden die Halbwertszeiten der Silberisotope $^{108}\mathrm{Ag}$ und $^{110}\mathrm{Ag}$ sowie des Indiumisotops $^{116}\mathrm{In}$ bestimmt werden. 

Die Untersuchung radioaktiven Zerfalls hat grundlegende Bedeutung für das Verständnis von Kernprozessen und ermöglicht Anwendungen in verschiedenen Bereichen der Physik, Medizin und Technik. Die Aktivierungsmethode stellt eine wichtige Technik dar, um stabile Isotope in radioaktive Quellen zu überführen, deren Eigenschaften dann präzise vermessen werden können. 

\section{Physikalische Grundlage}
\label{sec:PhyBasics}

\subsection{Erzeugung von Neutronenstrahlung}

Eine radioaktive Quelle kann durch Aktivierung stabiler Isotope mithilfe von Kernreaktionen hergestellt werden. Für die Anregung solcher Kernreaktionen werden häufig Neutronen verwendet, da diese elektrisch neutral sind und somit nicht der Coulomb-Wechselwirkung unterliegen. Der entsprechende Atomkern kann ein Neutron daher besonders leicht einfangen.

Für diesen Versuch wird eine Neutronenquelle eingesetzt, die aus einem Gemisch von Berylliumspänen und dem $\alpha$-Strahler Americium-241 ($^{241}\mathrm{Am}$) besteht. Es tritt folgende Kernreaktion auf:
\begin{equation}
    \label{eq:neutronenerzeugung}
    ^{9}\mathrm{Be} + \alpha \rightarrow {}^{12}\mathrm{C} + \mathrm{n}
\end{equation}

Bei dem Beschuss von Beryllium mit $\alpha$-Teilchen entsteht Kohlenstoff sowie ein freies Neutron. Das so erzeugte Neutron besitzt eine hohe kinetische Energie im Bereich von etwa $1$ bis $10\mathrm{MeV}$.

\subsection{Abbremsen der Neutronen (thermische Neutronen)}

Die bei der oben beschriebenen Reaktion entstehenden schnellen Neutronen müssen zunächst abgebremst werden, damit sie für bestimmte Kernreaktionen nutzbar gemacht werden können. Je langsamer ein Neutron ist, desto höher ist die Wahrscheinlichkeit, dass es von einem Atomkern eingefangen wird. Der Abbremsvorgang erfolgt hier durch einen Paraffinblock, der die Neutronenquelle umgibt.

Paraffin ist für diesen Prozess besonders geeignet, da es zahlreiche Wasserstoffkerne (Protonen) enthält. Trifft ein Neutron auf ein Proton, kommt es zu einem elastischen Stoß, bei dem das Neutron etwa die Hälfte seiner kinetischen Energie verliert. Bei Stößen mit Kohlenstoffkernen im Paraffin gehen dagegen weniger Energie verloren, da die Kohlenstoffkerne deutlich größer als die Neutronen sind. Die nach diesem Prozess abgebremsten Neutronen besitzen eine nahezu thermische Energie und werden als thermische Neutronen bezeichnet.

\subsection{Aktivierung von Silber (Erzeugung von Silberisotopen)}

Natürliches Silber setzt sich zusammen aus $51\%$ $^{107}\mathrm{Ag}$ und $49\%$ $^{109}\mathrm{Ag}$. Wird natürliches Silber wie oben erläutert aktiviert, entsteht aus $^{107}\mathrm{Ag}$ das Isotop $^{108}\mathrm{Ag}$ und aus $^{109}\mathrm{Ag}$ das Isotop $^{110}\mathrm{Ag}$. Beide entstandenen Isotope stellen $\beta$-Strahler dar, die durch Zerfall in $^{108}\mathrm{Cd}$ bzw. $^{110}\mathrm{Cd}$ übergehen.

Die Halbwertszeiten dieser beiden Silberisotope unterscheiden sich um etwa einen Faktor sechs. $^{110}\mathrm{Ag}$ weist eine kürzere Halbwertszeit auf und erreicht daher den Sättigungswert vor $^{108}\mathrm{Ag}$. Bei längerer Aktivierungszeit dominiert $^{108}\mathrm{Ag}$ aufgrund seiner deutlich längeren Halbwertszeit. Damit wird für kurze Aktivierungszeiten hauptsächlich $^{110}\mathrm{Ag}$ gebildet. Das Aktivitätsverhältnis der beiden Isotope verändert sich also mit der Bestrahlungszeit während der Aktivierung.

\img{Silber_plt}{Aktivierung der natürlichvorkommenden Silberisotope $51\%$ $^{107}\mathrm{Ag}$ und $49\%$ $^{109}\mathrm{Ag}$. \cite{skriptPAP22}}

Während des Aktivierungsprozesses werden kontinuierlich neue radioaktive Kerne gebildet, während gleichzeitig bereits vorhandene zerfallen. Die Aktivität $A(t)$ (Zahl der Zerfälle pro Sekunde) folgt dem Gesetz:
\eq{aktivitaet_buildup}{
    A(t) = A_{\infty}\left(1 - \mathrm{e}^{-\lambda t}\right)
}
Dabei bezeichnet $A_{\infty}$ die maximale Aktivität im Gleichgewicht (Sättigungsaktivität) und $\lambda$ beschreibt die Zerfallskonstante des Isotops. Für $t \rightarrow \infty$ nähert sich $A(t)$ asymptotisch an $A_{\infty}$ an.

\subsection{Aktivierung von Indium (Erzeugung von Indiumisotopen)}

Um stabiles Indium zu aktivieren, wird es analog bestreut. Durch die Bestrahlung kommen zur Bildung zweier Isomere eines radioaktiven Indiumisotops. Zum einen entsteht das Isotop $^{116}\mathrm{In}$, welches sich im energetischen Grundzustand befindet, zum anderen $^{116\mathrm{m}}\mathrm{In}$, welches sich im metastabilen Zustand befindet. Beide Isotope sind $\beta$-Strahler, zeigen aber unterschiedliche Halbwertszeiten. Beide Isotope des Indiums zerfallen zu $^{116}\mathrm{Sn}$ (Zinn).

Nach Abschluss der Aktivierung beginnt der rein radioaktive Zerfall, da keine neuen radioaktiven Kerne mehr erzeugt werden. Dann folgt:
\eq{aktivitaet_decay}{
    A(t) = A_0 \cdot \mathrm{e}^{-\lambda t}
}
Für die Halbwertszeit gilt somit:
\eq{halbwertszeit}{
    T_{1/2} = \frac{\ln(2)}{\lambda}
}
Zusätzlich kann auch die mittlere Lebensdauer $\tau$ definiert werden als:
\eq{lebensdauer}{
    \tau = \frac{1}{\lambda}
}

Durch Messung der Zerfallsrate über die Zeit lässt sich die Zerfallskonstante bestimmen, woraus sich gemäß \ceq{halbwertszeit} die Halbwertszeit berechnen lässt. Dies bildet die Grundlage der experimentellen Auswertung.

% \img{Nuklioide}{Ausschnitt der Nuklidkarte mit den Halbwertszeiten und der Energie der emittierten Strahlung\cite{skriptPAP22}}

\img[1]{NuDat_Nukleoide}{Ausschnitt Nuklidkarte nach NuDat (https://www.nndc.bnl.gov/nudat3/)}